\section{AI × Space 与新市场：从传感器到资源决策}

访谈的 space 讨论横跨 communication、Earth observation、navigation、农业、水与 lunar robotics。若只把这些项目并列，会变成宣传目录；更有教学价值的结构是 sensing-to-decision chain：space infrastructure 产生数据，AI 解释数据，政府或企业修改 physical action，再用结果校准模型。

\subsection{Communication、observation 与 harsh environment}

本节把课堂的 space 项目重新排列成一条 sensing-to-decision pipeline。读图时先区分 orbit 和 sensor 提供的原始能力，再追踪 imagery、positioning 与 connectivity 怎样进入 AI interpretation，最后检查 farming、water allocation 或 rover action 是否有地面反馈。只有最后一步闭环，space data 才从信息产品转化为资源治理能力。

\lecturefigure{14-ai-space-stack.png}{AI 与 space 的系统链从感知基础设施延伸到资源和物理决策。}{本地字幕 24:00--27:10、48:50--49:20；概念重绘。}

\begin{knowledgebox}{读图：Earth observation 的难点在 label 和 intervention}
卫星影像可以识别 land use、植被或水 stress，但模型输出只有进入 inspection、permit、irrigation 或 enforcement workflow 才改变结果。若地面 truth 稀缺、云层和季节变化大、政策干预没有记录，模型精度与真实节水效果之间会断裂。
\end{knowledgebox}

\begin{warningbox}{课堂中的节水和经济比例是讲者口径}
字幕给出 desertification、water saving、navigation 和 space-economy 的多个比例，但缺少统一方法和基线。本文不把这些数字作为定量结论，只保留“geospatial sensing + AI + policy action”需要闭环验证的系统模式。
\end{warningbox}

\subsection{NEOM 与 testbed：新市场不能免除证据要求}

新城市或大型项目可以减少 legacy integration，提供统一 identity、sensor 和 infrastructure 标准，因此适合 testbed；但“greenfield”也会放大 scope、capital、governance 和 adoption risk。系统验收仍要从小范围 service SLO、unit economics、safety incident 和 resident outcome 开始，而不是用项目愿景替代运行数据。

\begin{importantbox}{Testbed 的四级证据}
\begin{enumerate}
  \item prototype 能运行；
  \item pilot 在受控人群与环境中稳定；
  \item production 在真实负载、故障和组织约束下持续；
  \item transfer 能在不同城市、法规和供应商条件下复现。
\end{enumerate}
只有达到后两级，才说明平台能力具有可迁移性。
\end{importantbox}

\subsection{本章小结}

AI × space 的关键不在卫星或 rover 名单，而在 sensing、interpretation、decision 和 outcome 是否闭环。Greenfield testbed 可以加速实验，却不能降低对生产证据和迁移能力的要求。

